\section{Introduction}
Tabular data is the most ubiquitous data modality in enterprise machine learning, historically dominated by gradient-boosted decision trees (GBDTs) such as XGBoost \cite{chen2016xgboost} and LightGBM. Recently, the paradigm of In-Context Learning (ICL) has been successfully adapted from natural language processing to tabular data, culminating in Tabular Foundation Models (TFMs) like TabPFN \cite{hollmann2022tabpfn} and TabFM \cite{kong2026tabfm}. TabFM processes heterogeneous data directly via alternating row and column attention, entirely bypassing the gradient-based fine-tuning phase.

While extensive research has validated the efficiency and zero-shot scaling laws of TFMs, their security implications remain drastically underexplored. TFMs are uniquely vulnerable because their predictions on query data are strictly conditioned on the unverified, user-provided context window. If a malicious actor can manipulate the context window, they can hijack the model's predictive manifold.

In this paper, we expose a critical vulnerability: \textit{Adversarial Context Poisoning}. We demonstrate that imperceptible, schema-compliant perturbations to a small subset of the training context can dictate TabFM's predictions.

\textbf{Contributions:}
\begin{itemize}
    \item We introduce Continuous Subspace Poisoning, a mathematically rigorous Projected Gradient Descent (PGD) attack that optimizes continuous features while preserving non-differentiable categorical constraints.
    \item We demonstrate that the attack is highly transferable. Poisoned contexts generated by a Differentiable Kernel Regression surrogate successfully fool untouched TabFM instances in a black-box setting.
    \item We reveal through rigorous ablation studies that TabFM is hyper-fragile, suffering complete attention collapse even under a minuscule $1\%$ noise budget.
    \item We propose \textit{Context Subsampling}, a zero-cost, permutation-invariant defense mechanism that vastly outperforms traditional K-Nearest Neighbors sanitization, successfully shattering the adversarial gradients while improving baseline accuracy.
\end{itemize}
